\heading{量子力学A-24秋-期中2}
\noteinfo{整理自《量子力学 A》2024 年秋季学期 1 班期中考试参考答案；题面和解答按项目格式重写。}

\begin{question}
一维谐振子 $H=p^2/(2m)+m\omega^2x^2/2$。初态
\[
  \psi(x,0)=(Ax+Bx^3)\psi_0(x).
\]
\begin{enumerate}
  \item 求归一化条件。
  \item 测量能量的可能结果与概率。
  \item 求 $\psi(x,t)$。
  \item 求 $\expval{x},\expval{x^2},\expval{p},\expval{p^2}$ 并验证不确定关系。
\end{enumerate}
\begin{solution}
设 $s=\hbar/(2m\omega)$。由
\[
  x\ket0=\sqrt{s}\ket1,\qquad
  x^3\ket0=3s^{3/2}\ket1+\sqrt6\,s^{3/2}\ket3
\]
得
\[
  \ket{\psi(0)}=\sqrt s(A+3sB)\ket1+\sqrt6\,s^{3/2}B\ket3.
\]
归一化条件为
\[
  s|A+3sB|^2+6s^3|B|^2=1.
\]
测量能量只可能得到
\[
  E_1={3\over2}\hbar\omega,\qquad E_3={7\over2}\hbar\omega,
\]
概率分别为 $s|A+3sB|^2$ 与 $6s^3|B|^2$。时间演化
\[
  \ket{\psi(t)}=\sqrt s(A+3sB)\ee^{-3\ii\omega t/2}\ket1
  +\sqrt6\,s^{3/2}B\ee^{-7\ii\omega t/2}\ket3.
\]
态始终为奇宇称，故 $\expval{x}=\expval{p}=0$。记
\[
  c_1=\sqrt s(A+3sB),\qquad c_3=\sqrt6\,s^{3/2}B,
\]
其中 $|c_1|^2+|c_3|^2=1$。用
\[
  x^2=s(a^2+a^{\dagger2}+2a^\dagger a+1),\quad
  p^2={\hbar m\omega\over2}(-a^2-a^{\dagger2}+2a^\dagger a+1)
\]
得
\[
  \expval{x^2}=s\left(3|c_1|^2+7|c_3|^2
  +2\sqrt6\,\operatorname{Re}\!\left[c_1^*c_3\ee^{-2\ii\omega t}\right]\right),
\]
\[
  \expval{p^2}={\hbar m\omega\over2}\left(3|c_1|^2+7|c_3|^2
  -2\sqrt6\,\operatorname{Re}\!\left[c_1^*c_3\ee^{-2\ii\omega t}\right]\right).
\]
因此
\[
  \expval{x^2}\expval{p^2}\ge{\hbar^2\over4}.
\]
\end{solution}
\end{question}

\begin{question}
有限深方势阱与 $\delta$ 势叠加：
\[
V(x)=\begin{cases}
-V_0+V_1a\,\delta(x),& |x|<a,\\
0,& |x|>a,
\end{cases}
\qquad V_0,V_1,a>0.
\]
\begin{enumerate}
  \item 假设最低三个本征态为束缚态，定性画出波函数。
  \item 对束缚态 $E\in(-V_0,0)$，令 $k=\sqrt{2m(E+V_0)}/\hbar$，求 $k$ 方程。
  \item 求存在束缚态的参数条件。
\end{enumerate}
\begin{solution}
束缚态在 $|x|>a$ 指数衰减，阱内振荡；偶态在 $x=0$ 因排斥 $\delta$ 势出现朝向 $x$ 轴的尖点，奇态在 $x=0$ 为零且光滑。

记
\[
  \kappa=\sqrt{2mV_0/\hbar^2-k^2}.
\]
偶态可写为 $\psi=A\cos(k|x|-\phi)$（阱内），外部为 $B\ee^{-\kappa|x|}$。$x=0$ 的跳变条件给
\[
  \tan\phi={mV_1a\over\hbar^2k}.
\]
$x=a$ 处匹配给
\[
  k\tan(ka-\phi)=\kappa,
\]
即
\[
  {k\left[\tan(ka)-{mV_1a\over\hbar^2k}\right]
  \over1+\tan(ka){mV_1a\over\hbar^2k}}
  =\sqrt{2mV_0/\hbar^2-k^2}.
\]
奇态不受中心 $\delta$ 势影响：
\[
  -k\cot(ka)=\sqrt{2mV_0/\hbar^2-k^2}.
\]
至少存在一个偶束缚态的条件为
\[
  \sqrt{2mV_0\over\hbar^2}\,a>
  \arctan\left({mV_1a\over\hbar^2\sqrt{2mV_0/\hbar^2}}\right).
\]
\end{solution}
\end{question}

\begin{question}
三维简谐势
\[
  H={\vb p^2\over2m}+{m\omega^2\vb r^2\over2}.
\]
定义
\[
\varphi_1=\psi_{000},\quad \varphi_2=A_2x\psi_{000},\quad
\varphi_3=A_3y\psi_{000},\quad \varphi_4=A_4z\psi_{000}.
\]
\begin{enumerate}
  \item 求正的归一化因子 $A_2,A_3,A_4$。
  \item 求 $L_x,L_y,L_z$ 在该基底下的矩阵。
  \item 对态 $\frac12(\ket{\varphi_1}\ket{\downarrow}+\ket{\varphi_2}\ket{\uparrow}+\ket{\varphi_3}\ket{\uparrow}+\ket{\varphi_4}\ket{\uparrow})$，求测量 $J^2,J_z$ 的可能结果、概率与塌缩态。
\end{enumerate}
\begin{solution}
\[
  A_2=A_3=A_4=\sqrt{2m\omega\over\hbar}.
\]
角动量作用为
\[
  L_x\varphi_3=\ii\hbar\varphi_4,\quad L_x\varphi_4=-\ii\hbar\varphi_3,
\]
\[
  L_y\varphi_4=\ii\hbar\varphi_2,\quad L_y\varphi_2=-\ii\hbar\varphi_4,
\]
\[
  L_z\varphi_2=\ii\hbar\varphi_3,\quad L_z\varphi_3=-\ii\hbar\varphi_2,
\]
且 $L_a\varphi_1=0$。故三个矩阵只在 $\varphi_2,\varphi_3,\varphi_4$ 子空间中为 $\ell=1$ 表示。

轨道基底可取
\[
  \ket{1,0}=\ket{\varphi_4},\quad
  \ket{1,-1}={\ket{\varphi_2}-\ii\ket{\varphi_3}\over\sqrt2},\quad
  \ket{1,1}={-\ket{\varphi_2}-\ii\ket{\varphi_3}\over\sqrt2},
\]
以及 $\ket{0,0}=\ket{\varphi_1}$。再用 $\ell=0,1$ 与 $s=1/2$ 的 C-G 系数展开。非零概率和塌缩态为：
\[
\begin{array}{c|c|c}
(J^2,J_z)&P&\text{塌缩态}\\ \hline
(15\hbar^2/4,3\hbar/2)&1/4&
{1\over\sqrt2}(-\varphi_2-\ii\varphi_3)\ket{\uparrow}\\
(15\hbar^2/4,\hbar/2)&1/6&
\sqrt{2\over3}\varphi_4\ket{\uparrow}
 +{1\over\sqrt6}(-\varphi_2-\ii\varphi_3)\ket{\downarrow}\\
(15\hbar^2/4,-\hbar/2)&1/12&
\sqrt{2\over3}\varphi_4\ket{\downarrow}
 +{1\over\sqrt6}(\varphi_2-\ii\varphi_3)\ket{\uparrow}\\
(3\hbar^2/4,\hbar/2)&1/12&
-{1\over\sqrt3}\varphi_4\ket{\uparrow}
 +{1\over\sqrt3}(-\varphi_2-\ii\varphi_3)\ket{\downarrow}\\
(3\hbar^2/4,-\hbar/2)&5/12&
\sqrt{3\over5}\varphi_1\ket{\downarrow}
 +{1+\ii\over\sqrt{15}}(\varphi_2-\ii\varphi_3)\ket{\uparrow}
 -{1+\ii\over\sqrt{15}}\varphi_4\ket{\downarrow}
\end{array}
\]
\end{solution}
\end{question}

\begin{question}
三个自旋 $1/2$ 满足
\[
  H=J(\vb S_1\cdot\vb S_2+\vb S_2\cdot\vb S_3+\vb S_3\cdot\vb S_1)
  ={J\over2}\vb S^2-{9\over8}J\hbar^2,\quad J>0.
\]
初态为 $\ket{\psi(0)}=\ket{\downarrow\uparrow\uparrow}$。求 $\expval{S_{1z}}(t)$。
\begin{solution}
将初态分解为总自旋本征态：
\[
  \ket{\downarrow\uparrow\uparrow}
  ={1\over\sqrt3}\ket{3/2,1/2}
  +\sqrt{2\over3}\ket{1/2,1/2}.
\]
能量分别为 $E_{3/2}=3J\hbar^2/4$ 与 $E_{1/2}=-3J\hbar^2/4$。因此
\[
  \ket{\psi(t)}={1\over\sqrt3}\ee^{-\ii E_{3/2}t/\hbar}\ket{3/2,1/2}
  +\sqrt{2\over3}\ee^{-\ii E_{1/2}t/\hbar}\ket{1/2,1/2}.
\]
代回直积基并计算 $S_{1z}$，得
\[
  \expval{S_{1z}}(t)={\hbar\over2}\left[-{1\over9}
  -{8\over9}\cos\left({3J\hbar t\over2}\right)\right].
\]
\end{solution}
\end{question}
